\begin{study}{Backdoor, bars 9--12}
Handout note under \ch{IV}{}: \mel{3 \to \flat3}, ``backdoor'', \mel{\to 5}.
The 3rd of IV moves to $\flat$3 (IV becomes iv, minor), as in the
backdoor cadence back to I.
\end{study}

\begin{study}{Chromatic moments (main changes)}
Chords outside the key in \textcolor{chromatic}{red}; the melody note sounding
over them is circled, with its function on the chord:
\begin{itemize}\setlength{\itemsep}{0pt}
\item bar 13: \chx{\flat v}{\text{\o}7} \chx{VII}{7}, melody \mel{\tg{6}} = $\flat$7 of \chx{VII}{7} (V/iii);
\item bar 14: \chx{VI}{7} (V/ii), melody \mel{\tg{\hi{1}}} = $\sharp$9: tension against the major 3rd \mel{\sharp1} of the chord;
\item bar 15: \chx{II}{7} (V/V), melody \mel{\tg{7}} = 13;
\item bar 27: \chx{III}{7} (V/vi), melody \mel{\tg{\hi{2}}} = $\flat$7, resolving to \mel{\hi{1}} (3rd of \ch{vi}{-7}).
\end{itemize}
Characteristic altered degree of each chord: \chx{VII}{7}: \mel{\sharp2, \sharp4};
\chx{VI}{7}: \mel{\sharp1}; \chx{II}{7}: \mel{\sharp4}; \chx{III}{7}: \mel{\sharp5}.
\end{study}

\begin{study}{Ben Webster changes: secondary dominants}
Bars 1--2 and 17--18: \ch{I}{} \ch{IV}{(7)} $|$ \ch{iii}{-7} \ch{VI}{(7)} $|$ \ch{ii}{-7} \ch{V}{7}:
\ch{VI}{7} is V/ii. Bars 4 and 20: \ch{III}{7\flat9} is V/vi, leading to bar 5 (handout: ``expand!'').
Bars 8, 10, 24: \ch{I}{7} is V/IV. Bar 14: \ch{\flat III}{7} is a tritone substitute of \ch{VI}{7}.
Bar 28: \ch{\flat VI}{7} is a tritone substitute of \ch{II}{7}.
\end{study}

\begin{study}{Descending lines, bars 5--6, 21--22, 25--26}
Handout: ``descending'' under \ch{vi}{-7} and \ch{IV}{}. Bars 5--6 read
\ch{vi}{-7}, \ch{vi}{}/$\natural$7, \ch{vi}{}/7 (\ch{IV}{}): a line cliché with the bass moving down
under a sustained \ch{vi}{} chord. To be checked against the recording.
\end{study}
